\section{Expérience}
  \resumeSubHeadingListStart
    \resumeSubheading
      {Développeur Full-Stack et Principal}{2025 -- Présent}
      {Magicbox Logistics LLC}{À distance}
      \resumeItemListStart
        \resumeItem{Communique avec le gestionnaire pour créer des plans techniques et des documents de conception système}
        \resumeItem{Applique la modélisation statistique pour estimer les ventes de produits, la vitesse des tendances et le volume}
        \resumeItem{Conçoit et implémente des bases de données robustes pour stocker et suivre les produits tendance}
        \resumeItem{Utilise des outils d'automatisation de flux de travail pour optimiser l'efficacité du système et les opérations}
      \resumeItemListEnd

    \resumeSubheading
      {Développeur Full-Stack}{2024}
      {Boss Studios LLC}{À distance}
      \resumeItemListStart
        \resumeItem{Développé un jeu inspiré de Minecraft BedWars avec des mécaniques de jeu complètes}
        \resumeItem{Programmé le combat PvP, les systèmes de construction, les projectiles, les attaques au corps à corps et l'inventaire}
        \resumeItem{Intégré des systèmes de salons de jeu et des infrastructures multijoueurs}
      \resumeItemListEnd

    \resumeSubheading
      {Développeur Frontend}{2024}
      {Oakley Productions LLC}{À distance}
      \resumeItemListStart
        \resumeItem{Contribué à « Anime Defenders », un jeu de Tower Defense sur Roblox (+3,4 milliards de visites, pic de 100K connexions)}
        \resumeItem{Développé la logique de jeu pour plus de 10 unités/personnages et 30 capacités au total}
        \resumeItem{Conçu des effets visuels évolutifs en utilisant des systèmes de particules, des trajectoires de projectiles et des courbes de Bézier}
      \resumeItemListEnd

    \resumeSubheading
      {Développeur de Jeux Full-Stack}{2020 -- Présent}
      {Freelance/Contractuel}{À distance}
      \resumeItemListStart
        \resumeItem{Appliqué les mathématiques (matrices, vecteurs, cinématique inverse, calculs d'intersections) aux mécaniques de jeu}
        \resumeItem{Appliqué la physique (cinématique, mouvement de projectiles, ressorts, réflexion de lumière) pour des simulations réalistes}
        \resumeItem{Développé des systèmes réseau, animations personnalisées, effets visuels, UI/HUD, pathfinding IA et bases de données}
      \resumeItemListEnd
  \resumeSubHeadingListEnd

\section{Projets}
  \resumeSubHeadingListStart
    \resumeProjectHeading
      {\textbf{HackTheMountain Hackathon (vibe-fm)} $|$ \emph{Python, TypeScript, Expo, React Native, OpenRouter}}{2026}
      \resumeItemListStart
        \resumeItem{Conçu, développé et présenté l'application mobile en solo à Polytechnique Montréal}
        \resumeItem{Intégré l'API Shazam via RapidAPI pour l'identification musicale en temps réel}
        \resumeItem{Utilisé les bibliothèques natives Expo pour l'enregistrement audio et les vibrations haptiques}
        \resumeItem{Analysé les pistes audio à l'aide de FreqBlog (BPM, énergie, valence, ambiance, etc.)}
        \resumeItem{Exploité Gemini Flash via OpenRouter pour synthétiser les données qualitatives et profiler les utilisateurs}
      \resumeItemListEnd

    \resumeProjectHeading
      {\textbf{levlrai} $|$ \emph{TypeScript, Next.js, Supabase, Stripe, Redis, Vite, Tailwind}}{2026}
      \resumeItemListStart
        \resumeItem{Architecturé une plateforme d'étude IA avec parsing de cours via Mammoth et PDF-Parse}
        \resumeItem{Intégré Supabase Postgres, Google OAuth, limitation de débit via Upstash Redis et abonnements Stripe}
        \resumeItem{Intégré l'API Stripe pour les paiements de bout en bout par abonnement}
        \resumeItem{Maintenu une haute fiabilité du code grâce à des tests CI utilisant Playwright et Vitest}
      \resumeItemListEnd

    \resumeProjectHeading
      {\textbf{Simulation CFD} $|$ \emph{Simulation Physique, Moteur de Jeu, C\#}}{2026}
      \resumeItemListStart
        \resumeItem{Programmé une simulation de dynamique des fluides (CFD) sous les contraintes d'un moteur de jeu}
        \resumeItem{Implémenté un algorithme de pilotage inspiré des boids et des forces de pression pour simuler des trajectoires de particules}
        \resumeItem{Appliqué le principe de Bernoulli pour modéliser les relations de vitesse et de pression en temps réel}
      \resumeItemListEnd

    \resumeProjectHeading
      {\textbf{Serveur Market MCP (AI Hackfest Hackathon)} $|$ \emph{Python, FastMCP, BeautifulSoup4}}{2026}
      \resumeItemListStart
        \resumeItem{Créé un serveur Model Context Protocol (MCP) d'intelligence de marché pour permettre aux agents d'IA d'interroger des données en temps réel issues des wikis de l'écosystème Roblox, Rolimons et des API Cloud}
        \resumeItem{Utilisé FastMCP et BeautifulSoup4 pour concevoir des flux de travail d'extraction et de collecte de données de haute fiabilité}
      \resumeItemListEnd

    \resumeProjectHeading
      {\textbf{Job Scout (Testsprite Hackathon)} $|$ \emph{React, Vite, FastAPI, Python, Vercel, Railway}}{2026}
      \resumeItemListStart
        \resumeItem{Dirigé une équipe de développement de trois personnes pour créer un outil de correspondance d'emplois par IA utilisant l'analyse de CV par Gemini, la récupération d'emplois par SerpAPI et la génération automatique de lettres de motivation}
        \resumeItem{Déployé l'interface utilisateur sur Vercel avec React et Vite, couplée à un backend FastAPI hébergé sur Railway}
      \resumeItemListEnd
  \resumeSubHeadingListEnd

\section{Activités Parascolaires}
  \resumeSubHeadingListStart
    \resumeSubheading
      {Mentor en Développement de Jeux \& Programmation}{2023 -- Présent}
      {Mentorat entre pairs / Tutorat}{À distance}
      \resumeItemListStart
        \resumeItem{Guidé et tutoré plusieurs pairs dans des cours de programmation informatique}
        \resumeItem{Mentoré un étudiant de fin de cycle pour son projet d'intégration Unity en 200.C1 (2025)}
        \resumeItem{Enseigné des concepts de développement de jeux : cinématique inverse FABRIK, pathfinding A*, mouvement de projectiles, logique de collision}
      \resumeItemListEnd

    \resumeSubheading
      {Bénévole aux Portes Ouvertes}{Coll\`{e}ge Lionel-Groulx}{Sainte-Th\'{e}r\`{e}se, QC}
      {2025}
      \resumeItemListStart
        \resumeItem{Accueilli et guidé les visiteurs dans le programme de sciences pures et appliquées sur le campus}
        \resumeItem{Répondu aux questions sur les cours, la charge de travail académique et les attentes}
      \resumeItemListEnd
  \resumeSubHeadingListEnd

\section{Compétences Techniques}
  \begin{itemize}[leftmargin=0.15in, label={}]
    \small{\item{
      \textbf{Langages}{: C\#, Python, JavaScript, TypeScript, Lua, SQL (Postgres + SQLite), HTML/CSS} \\
      \textbf{Frameworks}{: React, FastAPI, Next.js, Playwright, Vitest, React Native, Tailwind} \\
      \textbf{Outils de Développement}{: Git, GitHub, Vercel, Railway, Expo, Supabase, Redis, VS Code, Visual Studio, Vite, Fusion360, Unity, Blender, Roblox Studio} \\
      \textbf{Outils d'IA}{: Claude Code, Antigravity, Hermes, Gemini CLI, OpenCode, Goose} \\
      \textbf{Bibliothèques}{: PyTrends, NumPy, Matplotlib, Mammoth, PDF Parse}
    }}
  \end{itemize}
